\documentclass[10pt,a4paper]{amsart}
\usepackage[margin=19mm]{geometry}
\usepackage{amsmath,amssymb,mathtools}
\usepackage[scr=rsfs,cal=euler]{mathalfa}
\usepackage{xfrac}
\usepackage[unicode,hidelinks,bookmarks=false]{hyperref}
\newcommand{\ie}{i.e.\ }
\newcommand{\cf}{cf.\ }
\newcommand{\resp}{respectively\ }
\newcommand{\Subsection}{\subsection}
\providecommand{\nobreakdash}{\nobreak\hbox{-}}
\setlength{\emergencystretch}{2em}
\multlinegap0pt
\usepackage{amssymb}
\usepackage{bm}
\usepackage{mathtools}
\usepackage[shortlabels]{enumitem}
\usepackage{xcolor}
\usepackage{tikz}
\usepackage{dsfont}
\usepackage{float}
\newtheorem{lemma}{Lemma}
\newtheorem{claim}{Claim}
\newcommand\eps{\varepsilon}
\title{Robustness and hyperstability for the Erd\H{o}s--Gallai theorem}
\author{Micha Christoph, Alp M\"uyesser, Yuval Wigderson}
\date{Source: arXiv:2607.02483v2 (CC BY 4.0)}
\begin{document}
\maketitle

\noindent\emph{Source and licence.} Robustness and hyperstability for the Erd\H{o}s--Gallai theorem. Version arXiv:2607.02483v2 (CC BY 4.0), distributed by arXiv under the Creative Commons Attribution 4.0 International licence, http://creativecommons.org/licenses/by/4.0/, as recorded in the arXiv licence metadata for this exact version and shown on the abstract page https://arxiv.org/abs/2607.02483v2. Full licence notice: \texttt{licenses/arxiv-2607.02483-CC-BY-4.0.txt}.\par\smallskip

\noindent\textit{Editorial scope.} Counted unit: the article's lemma regularity (source0.tex lines 299--305). The article's complete original proof of it is delivered with it (source0.tex lines 417--464, one \texttt{proof} environment of 48 lines, which the article places away from the statement). No mathematics is written, repaired or re-derived: the statement and the proof are the article's own lines, in the article's own notation, and no retained line is substituted or re-flowed.  Retained uncounted as the source-local context the proof needs: lines 312--314; lines 319--321; lines 325--327; lines 408--409.  Nothing was omitted from any retained span, so the package carries no omission record.  No retained line contains a citation, so the wrapper carries no bibliography and no bibliography entry.  No retained line contains a figure environment, an image-inclusion command or drawing code, so no image is delivered and the package declares no asset dependency.  Editorial formatting only, in addition: an AMS \texttt{amsart} preamble that restates the article's own declarations and macros used by the retained lines, namely the environments \texttt{lemma}, \texttt{claim} on separate counters, together with the standard packages those lines need.  The retained environments print in this document as Lemma 1, Claim 1 in order of appearance, whereas the article prints the same items with the numbering of its own full document.  The complete native source of the article as distributed in the arXiv e-print for this exact version is bundled unchanged as \texttt{sources/204318/source0.tex}.
\begin{lemma}\label{lem: random joined}
    For every $\varepsilon>0$ the following holds for large enough $k$ and $t\geq k$. Let $(A,B)$ be an $\eps$-regular pair with $|A|,|B|=t$ of density at least $2\varepsilon$. For $p\geq k/t$, with probability at least $1-e^{-t}$, $(A,B)_p$ is $\varepsilon t$-joined.
\end{lemma}

\begin{lemma}\label{lem: DFS bipartite holes}
    Let $G=(A,B)$ be an $m$-joined, balanced bipartite graph. Then there exists a path in $G$ of length $|V(G)|-4m$.
\end{lemma}

\begin{lemma}\label{lem: regularity connector}
    Let $G$ be a graph on vertex set $U_1\cup\ldots\cup U_r$ such that for every $1\leq i< r$ the pair $(U_i,U_{i+1})$ is $m$-joined. Suppose further that each $U_i$ is of size at least $3m$ and let $Z\subseteq V(G)$ be of size at most $m$. Then there exists a path in $G$ from $U_1$ to $U_r$ avoiding $Z$ such that the $i$th vertex is in $U_i$ for all $i \in [r]$.
\end{lemma}

\begin{lemma}\label{lem:cut-dense case} For every $C,M,\rho,q,\xi>0$ the following holds for large enough $K$ and every $d\geq K$. Let $H$ be a $\rho$-cut-dense graph of order at most $Md$ which contains a $qd$-regular subgraph of order at least $Cd$. For $p\geq K/d$, $H_p$ contains a cycle of length $(1-\xi)Cd$ with probability at least $1-e^{-\Omega(d)}$.
\end{lemma}

\begin{lemma}[Szemer\'edi's regularity lemma]\label{lem: regularity}
    For every $\varepsilon,p>0$ and $L_0\in\mathbb N$ there exists $L\geq L_0$ such that the following holds. Let $G,H$ be two graphs on the same vertex set $V$ and let $U\subseteq V$ with $|U|\geq p|V|$. Then, there exists an equipartition $V=V_0\sqcup\ldots\sqcup V_m$ for some $L_0\leq m\leq L$ such that
    \begin{itemize} 
        \item for every $1\leq i\leq m$, $V_i$ is either contained in $U$ or disjoint from $U$, and 
        \item for every $1\leq i\leq m$ for all but at most $\varepsilon m$ choices of $1\leq j\leq m$ with $i\neq j$ we have that $(V_i,V_j)$ is $\varepsilon$-regular in $G$ and $H$.
    \end{itemize}
\end{lemma}

\begin{proof}[Proof of Lemma~\ref{lem:cut-dense case}]
We fix constants satisfying the hierarchy: $1/C,1/M,\rho,q,\xi \gg \varepsilon \gg L_0^{-1} \gg K^{-1} \gg d^{-1}$. Let $G\subseteq H$ be a $qd$-regular subgraph of order at least $Cd$.

Let $V_0,\ldots,V_m$ be obtained from Lemma~\ref{lem: regularity} applied to $G$ and $H$ with $\varepsilon, p=C/M\leq |V(G)|/|V(H)|, L_0$ and $U=V(G)$. By relabeling if need be, let us assume that $V(G)\setminus V_0=V_1\sqcup\ldots\sqcup V_{m_G}$. For simplicity of presentation, we assume henceforth that each $V_i$ has size exactly $t=|V(H)|/(m+1)$. For every $1\leq i\leq m$, fix an arbitrary subset $V_i'\subseteq V_i$ of size $3\varepsilon t$.

The following is the only use we make of randomness. By Lemma~\ref{lem: random joined} and a union bound, with probability at least $1-2m^2e^{-t}=1-e^{-\Omega(d)}$ for each $1\leq i<j\leq m$, if $(V_i,V_j)$ is $\varepsilon$-regular with density at least $2\varepsilon$ in either $H$ or $G$, then it is $\varepsilon t$-joined in $H_p$. For the remaining argument, let us fix an outcome of the randomness where this holds.

For each $1\leq i<j\leq m_G$, set $\mu(V_i,V_j)=\left\lfloor\frac{e_G(V_i,V_j)}{2\varepsilon t^2}\right\rfloor$ whenever $(V_i,V_j)$ is $\varepsilon$-regular in $G$, otherwise set $\mu(V_i,V_j)=0$. Let $\mathcal R_G$ be the multigraph on $\{V_1,\ldots,V_{m_G}\}$ where we add exactly $\mu(V_i,V_j)$ edges between every pair $V_i,V_j$ with $1\leq i<j\leq m_G$. Next, we aim to show that $\mathcal R_G$ is almost regular so that we can apply the Vizing--Gupta theorem to conclude that $\mathcal R_G$ contains an almost spanning matching. For $1\leq i\leq m_G$, $V_i$ is incident to at most $t\cdot qd$ edges in $G$. Note also that $V_i$ is incident to at least $\deg_{\mathcal{R}_G}(V_i)\cdot 2\varepsilon t^2$ edges in $G$. We get that $\deg_{\mathcal{R}_G}(V_i)\leq qd/(2\varepsilon t)$. As this holds for all $1\leq i\leq m_G$, it is also an upper bound on $\Delta(\mathcal{R}_G)$. On the other hand, for every $1\leq i<j\leq m_G$, we have that $e_G(V_i,V_j)\leq(\mu(V_i,V_j)+1)\cdot 2\varepsilon t^2$ whenever $(V_i,V_j)$ is $\varepsilon$-regular in $G$ and otherwise $e_G(V_i,V_j)\leq t^2$. Furthermore, for every $1\leq i\leq m_G$, we have $e_G(V_i)\leq \binom{t}{2}$ and $e_G(V_0,V(G)\setminus V_0)\leq t\cdot qd$.

Therefore, we get
\begin{align*}
    qd|V(G)|/2=e(G)\leq e(\mathcal{R}_G)\cdot 2\varepsilon t^2+\binom{m_G}{2}\cdot 2\varepsilon t^2+m_G\cdot\varepsilon m\cdot t^2+m_G\cdot \binom{t}{2}+t\cdot qd.
\end{align*}
Recall that $|V(G)|\geq m_G\cdot t$ and that $(m+1)t=|V(H)| \leq Md$, and hence $d/t\geq m/M$. Dividing by $2\varepsilon t^2$ and rearranging now yields
\begin{align*}
    e(\mathcal R_G) \geq qd/(4\varepsilon t)\cdot m_G- \binom{m_G}{2}-m_Gm/2-m_G/(4\varepsilon)-qd/(2\varepsilon t)\geq (1-5M\varepsilon/q-2/m_G)qd/(4\varepsilon t)\cdot m_G. 
\end{align*}
Using that $1/m_G\ll M\varepsilon/q$, we get 
$$
e(\mathcal R_G) \geq (1-6M\varepsilon/q)qd/(4\varepsilon t)\cdot m_G. 
$$
Finally, denote by $\mu(\mathcal R_G)$ the maximum edge multiplicity of $\mathcal R_G$. Then $\mu(\mathcal R_G)\leq 1/(2\varepsilon)$. By the Vizing--Gupta theorem, we have that $\mathcal R_G$ has a proper edge-coloring with $\Delta(\mathcal R_G)+\mu(\mathcal R_G)\leq (qd+t)/(2\varepsilon t)$ colors. By averaging, one of those colors contains at least
$$
\frac{e(\mathcal R_G)}{(qd+t)/(2\varepsilon t)}\geq (1-6 M \varepsilon /q)m_G/2\cdot\frac{qd}{qd+t}\geq (1-7M\varepsilon/q)m_G/2
$$
edges. Let $F\subseteq \mathcal R_G$ be the matching corresponding to this color and note that $|V(F)|\geq (1-7M\varepsilon/q) m_G$.

For $V_iV_j\in E(F)$, we have that $(V_i\setminus V_i',V_j\setminus V_j')$ is $\varepsilon t$-joined in $H_p$, since $(V_i,V_j)$ is $\varepsilon t$-joined by assumption and, hence, every induced subgraph of $(V_i,V_j)$ is $\varepsilon t$-joined. By Lemma~\ref{lem: DFS bipartite holes}, there is a path of length $|V_i\setminus V_i'|+|V_j\setminus V_j'|-4\varepsilon t = 2t-10\varepsilon t$ in $H_p$.  Let $P_1,\ldots,P_{|F|}$ be such paths for each edge in $F$. We aim to connect these paths into a cycle. For every $P_i$, give it an arbitrary direction and let $I_i$ be the initial segment of $6\varepsilon t$ vertices, $O_i$ the segment of the last $6\varepsilon t$ vertices and $P_i'$ the middle segment which is of length at least $2t-22\varepsilon t$. For every $1\leq i\leq |F|$, let $W_i$ denote one of the two clusters which intersect $P_i$.

Let $\mathcal R_H$ be the graph on $\{V_1,\ldots, V_m\}$ where $V_i,V_j$ are connected if $(V_i,V_j)$ is $\varepsilon$-regular of density at least $2\varepsilon$ in $H$.
    \begin{claim}\label{claim: cluster graph is connected}
        $\mathcal R_H$ is connected.
    \end{claim}
    \begin{proof}
        Suppose otherwise, without loss of generality, let us say that $\mathcal R_H$ does not contain an edge between $\{V_1,\ldots, V_\ell\}$ ($\ell\leq m/2$) and $\{V_{\ell+1},\ldots,V_m\}$. Then, all edges in $H$ between $A=\bigcup_{i=1}^\ell V_i$ and $V(H)\setminus A$ are between pairs $(V_i,V_j)$ with either $1\leq i\leq\ell$ and $j=0$, or, $1\leq i\leq\ell<j\leq m$ and either $(V_i,V_j)$ is not $\varepsilon$-regular or $(V_i,V_j)$ has density less than $2\varepsilon$. The number of such edges is at most
        $$
        \ell t^2+\ell\varepsilon m\cdot t^2+\ell(m-\ell)\cdot2\varepsilon t^2\leq \ell(m+1-\ell)t^2\cdot 4\varepsilon< \rho\ell(m+1-\ell)t^2,
        $$
        using that $\varepsilon \ll \rho$.
        Noting that $\ell t=|A|$ and $(m+1-\ell)t=|V(H)\setminus A|$ this contradicts the assumption that $H$ is $\rho$-cut-dense.
    \end{proof}
    We iteratively find disjoint paths $Q_i$ in $H_p$ for every $1\leq i\leq |F|$ of length at most $m$ each such that $Q_i$ connects $O_i$ with $I_{i+1}$ (indices mod $|F|$) and the internal vertices are in $\bigcup_{j=1}^m V_{j}'$. Suppose we have already found $Q_1,\ldots, Q_{i-1}$. Set $Z=\bigcup_{j=1}^{i-1}V(Q_{j})$ and note that $|Z|\leq m^2\leq \varepsilon t$, using that $d^{-1}\ll \varepsilon$. Let $P=(V_{j_1},\ldots, V_{j_r})$ be a path in $\mathcal R_H$ with $V_{j_1}=W_i$ and $V_{j_r}=W_{i+1}$. Since $\mathcal R_H$ is a connected $m$-vertex graph, such a path exists and is of length at most $m$. Since $V_{j_h}V_{j_{h+1}}\in \mathcal R_H$, the pair $(V_{j_h},V_{j_{h+1}})$ is $\varepsilon t$-joined in $H_p$. Note that $|O_i\cap W_i|\geq 3\varepsilon t$ and also $|I_{i+1}\cap W_{i+1}|\geq 3\varepsilon t$. By Lemma~\ref{lem: regularity connector} applied with $U_1=O_i\cap W_i$, $U_2 = V_{j_2}',\ldots, U_{r-1}=V_{j_{r-1}}'$ and $U_r=I_{i+1}\cap W_{i+1}$, there exists a desired path $Q_i$ in $H_p$. 

    Finally, concatenating all the pieces, we conclude that there is a cycle $C$ in $H_p$ which includes $P_i'$ for all $1\leq i\leq |F|$.
    Thus, the total length of this cycle is at least 
    $$\sum_{i=1}^{|F|}|P_i'|\geq (1-7M\varepsilon/q)m_G/2\cdot (2-22\varepsilon)t\geq (1-7M\varepsilon/q)(1-12\varepsilon)|V(G)|\geq (1-\xi)Cd,$$
    using that $\varepsilon \ll \xi,q,M^{-1}$. This concludes the proof.
\end{proof}

\end{document}
